\begin{tikzpicture}[x=1mm,y=1mm,font=\figurefont\footnotesize,text=NNInk,
  nn flow/.append style={draw=NNWire,rounded corners=1.5mm},
  nn operator path/.append style={draw=NNWire,rounded corners=1.5mm},
  nn operator/.append style={fill=NNHidden!18}]
  \path[use as bounding box] (0,-60) rectangle (112,43);
  \node[anchor=east] (old) at (10,36) {$\bm h^{(t-1)}$};
  \node[anchor=east] (x) at (10,-28) {$\bm x_t$};
  \node[anchor=west] (new) at (106,19) {$\bm h^{(t)}$};
  \node[anchor=west] (y) at (106,-48) {$\widehat{\bm y}_t$};
  \node[nn operator,draw=NNInput,fill=NNInput!18,
    minimum width=16mm,minimum height=13mm] (in) at (27,-28)
    {组合输入};
  \node[nn operator,minimum width=22mm,minimum height=13mm]
    (gates) at (56,-28) {Sigmoid门\\$\bm r^{(t)},\bm z^{(t)}$};
  \node[nn operator,minimum width=24mm,minimum height=15mm]
    (cand) at (56,-2) {重置＋候选\\$\widetilde{\bm h}^{(t)}$};
  \node[nn pointwise] (keep) at (78,30) {$\odot$};
  \node[nn pointwise] (write) at (78,8) {$\odot$};
  \node[nn pointwise] (add) at (94,19) {$+$};
  \node[nn operator,draw=NNOutput,fill=NNOutput!18,minimum width=18mm]
    (readout) at (94,-48) {读出$g$};

  % 旧状态同时进入保留路径与组合输入。
  \draw[nn operator path] (old.east)--(69,36)--(69,30)--(keep.west);
  \draw[nn operator path] (27,36)--(in.north);
  \fill[NNWire] (27,36) circle[radius=.45mm];
  \draw[nn flow] (x.east)--(in.west);
  \draw[nn operator path] (in.east)--(gates.west);
  \draw[nn operator path] (40,-28)--(40,-2)--(cand.west);
  \fill[NNWire] (40,-28) circle[radius=.45mm];

  % 重置门进入候选；更新门在空白路由带中分成保留与写入两路。
  \draw[nn operator path] (gates.north)--node[right] {$\bm r^{(t)}$} (cand.south);
  \draw[nn operator path] (gates.east)--(70,-28)--(70,19);
  \draw[nn operator path] (70,19)--(70,26)--(78,26)--(keep.south);
  \draw[nn operator path] (70,19)--(70,4)--(78,4)--(write.south);
  \fill[NNWire] (70,19) circle[radius=.45mm];
  \node[anchor=east] at (68,24) {$\bm z^{(t)}$};
  \node[anchor=east] at (68,11) {$\bm1-\bm z^{(t)}$};
  \draw[nn operator path] (cand.north)--(57,8)--(write.west);
  \draw[nn operator path] (keep.east)--(86,30)--(86,23)--(add.north west);
  \draw[nn operator path] (write.east)--(86,8)--(86,15)--(add.south west);
  \draw[nn operator path] (add.east)--(new.west);
  \draw[nn flow] (add.south)--(readout.north);
  \draw[nn flow] (readout.east)--(y.west);
\end{tikzpicture}
